%!TEX program = xelatex
\documentclass[11pt, a4paper]{article}

% ─── Packages ────────────────────────────────────────────────────────────────
\usepackage{fontspec}
\usepackage{kotex}
\usepackage[margin=2.5cm]{geometry}
\usepackage{amsmath, amssymb, bm}
\usepackage{graphicx}
\usepackage{booktabs}
\usepackage{multirow}
\usepackage{array}
\usepackage{xcolor}
\usepackage{hyperref}
\usepackage{natbib}
\usepackage{caption}
\usepackage{subcaption}
\usepackage{float}
\usepackage{algorithm}
\usepackage{algpseudocode}
\usepackage{enumitem}
\usepackage{parskip}
\usepackage{titlesec}
\usepackage{setspace}

\setmainhangulfont{Malgun Gothic}
\setsanshangulfont{Malgun Gothic}
\setmonohangulfont{Malgun Gothic}

\onehalfspacing
\setlength{\parindent}{0pt}
\setlength{\parskip}{6pt}

\hypersetup{
    colorlinks=true,
    linkcolor=blue!60!black,
    citecolor=green!50!black,
    urlcolor=blue!70!black
}

\graphicspath{{figures/}{../exp032c_figures/}{../exp033_figures/}{../exp034_figures/}{../exp035_figures/}{../phase15_figures/}{../phase16d_figures/}}

% ─── Title ───────────────────────────────────────────────────────────────────
\title{
    \vspace{-1cm}
    {\large 캡스톤 디자인 프로젝트 — 최종 보고서}\\[0.8em]
    \textbf{BTC 강화학습 그리드 트레이딩에서 Reward 함수 설계의 영향:\\
    Pareto 프론티어, Winner Reversal,\\
    그리고 Prospect Theory의 분포 외 강건성}
}

\author{
    이찬희 (Lee Chanhee)\\
    컴퓨터공학부\\
    서울과학기술대학교\\
    \texttt{happilyfly@seoultech.ac.kr}
}

\date{2026년 6월}

% ─── Document ────────────────────────────────────────────────────────────────
\begin{document}

\maketitle
\thispagestyle{empty}

% ─── Abstract (한글) ─────────────────────────────────────────────────────────
\begin{abstract}
본 논문은 BTC/USDT 1시간봉 시장에서 가격 방향 예측 없이 ATR(Average True Range) 비례
지정가 그리드를 결정하는 PPO 기반 강화학습 정책에 대해 \emph{reward 함수 설계가 정책의
행동 패턴과 일반화 성능에 미치는 영향}을 정량 검증한다. 4가지 reward 변형 ---
대칭 보상(\texttt{sym}), 비대칭 보상(\texttt{asym}, $\beta=3.42$),
미분 샤프 비율(\texttt{dsr}, $\eta=1/28\,\mathrm{h}$),
프로스펙트 이론 보상(\texttt{pt}, $\alpha=0.68, \lambda=3.30$) --- 을 동일 환경에서
각각 10 시드 $\times$ 100만 스텝 학습 후, (i) 단일 분할 검증 (Val 2021--2023),
(ii) 6중 CPCV(Combinatorial Purged CV) 15-path 평가,
(iii) 봉인 해제된 테스트 셋(Test 2024--2026) 분포 외 평가의 세 환경에서 비교한다.

주요 발견:
(1) Val에서 4개 변형이 ATR 기반 베이스라인(Val Sharpe $\approx$ 1.38--1.51)을 모두
$+11\sim24\%$ 초과하나 단일 winner는 없으며, 두 클러스터
\{\texttt{sym}, \texttt{dsr}\}(공격적) vs \{\texttt{asym}, \texttt{pt}\}(보수적)로
통계적으로 분리된다(그룹 내 $|d|<0.3$, 그룹 간 $|d|>0.8$).
(2) Reward 형식이 정책의 거래 빈도와 holding 시간을 직접 결정하며
(정책 거리 비율 2.22배), \texttt{dsr}의 hold rate가 다른 변형의 2--6배인 것이
sliding-window reward 자체의 메모리 구조 결과임을 정량 확인한다.
(3) 평가 환경별로 winner가 reversal: Val의 \texttt{sym}(1.87) $\to$
CPCV의 \texttt{dsr}(1.41, p$<$0.001) $\to$ Test의 \texttt{pt}(0.34--0.37,
두 소스 일관 p$<$0.002).
(4) \texttt{pt}의 분포 외 강건성 메커니즘: $\lambda=3.30$ 손실 회피와 $\alpha=0.68$
오목한 효용함수가 정책의 매수 후 평균 1.4시간(최대 6시간) 빠른 청산을 학습시켜
미관측 강세장(BTC \$42K $\to$ \$76K)의 매도 타이밍 위험을 회피한다. 반면 \texttt{dsr}은
평균 4.58시간(최대 169시간 = 7일) holding을 학습시켜 동일 환경에서 최악 성능 기록.
(5) Val--Test 분포 차이가 통계적으로 매우 유의(KS $p<10^{-10}$, 변동성 $-27\%$)하며,
모든 정책의 Val/Test 행동은 거의 동일($\Delta$ action $\leq 5\%$)이므로 OOS 성능 차이는
정책이 아닌 \emph{시장 분포 이동}이 직접 원인이다.

본 연구의 학술적 기여는 reward 설계의 효과가 단일 Sharpe metric의 ``alpha source''가
아니라 \emph{risk profile 차원의 trade-off + 평가 환경별 winner reversal}로 나타남을
정량 보이는 것이며, Kahneman과 Tversky(1979)의 prospect theory가 강화학습 정책에 부여하는
분포 외 안전성을 첫 정량 확인한 것이다. 코드와 데이터:
\url{https://github.com/cosmicpotato2047/capstone-rl-trading}.
\end{abstract}

\vspace{1em}
\noindent\textbf{키워드:} 강화학습, 그리드 트레이딩, PPO, Reward Design, Prospect Theory,
Differential Sharpe Ratio, CPCV, 분포 외 일반화, 암호화폐

% ─── Abstract (영문) ─────────────────────────────────────────────────────────
\vspace{1em}
\begin{center}\textbf{Abstract (English)}\end{center}
\small
We quantitatively study the effect of \emph{reward function design} on PPO-based RL
policies for ATR-proportional grid trading on the BTC/USDT 1-hour market. Four reward
variants---symmetric (\texttt{sym}), asymmetric ($\beta=3.42$, \texttt{asym}),
Differential Sharpe Ratio ($\eta=1/28\,\mathrm{h}$, \texttt{dsr}), and
prospect-theoretic ($\alpha=0.68$, $\lambda=3.30$, \texttt{pt})---are each trained
with 10 seeds $\times$ 1M steps and evaluated on three environments:
(i) single-split validation (Val 2021--2023), (ii) 6-fold CPCV 15 paths, and
(iii) the unsealed out-of-sample Test (2024--2026, BTC \$42K\,$\to$\,\$76K).

Key findings:
(1) All four variants exceed the ATR baseline by $+11$--$24\%$ on Val, but partition
into two clusters \{\texttt{sym}, \texttt{dsr}\} (aggressive) vs.
\{\texttt{asym}, \texttt{pt}\} (conservative); within-cluster $|d|<0.3$ vs.\
across-cluster $|d|>0.8$.
(2) Mechanism: reward asymmetry directly determines trade frequency and holding time
(policy distance ratio 2.22$\times$), and the sliding-window memory of DSR results in
2--6$\times$ higher hold rate.
(3) Winner reverses across evaluation environments: Val=\texttt{sym}(1.87)
$\to$ CPCV=\texttt{dsr}(1.41, p$<$0.001) $\to$ Test=\texttt{pt}(0.34--0.37,
consistent across two sources, p$<$0.002).
(4) Mechanism for \texttt{pt}'s OOS robustness: short holding (mean 1.4\,h, max 6\,h)
avoids sell-side timing risk in the unseen bull market. In contrast, \texttt{dsr}
learns long holding (mean 4.58\,h, max 169\,h = 7 days) and records the worst OOS
performance, showing that the same reward formulation drives both in-sample advantage
and OOS failure as two sides of the same coin.
(5) Val--Test distribution shift is statistically significant
(KS $p<10^{-10}$, $-27\%$ variance), while policy actions remain nearly invariant
($\Delta \leq 5\%$) between Val and Test, indicating the regime shift, not the policy,
drives OOS differences. Our main contribution is the quantitative demonstration that
reward design exerts its effect through the trade-off dimension rather than a single
Sharpe ``alpha source,'' and provides the first quantitative evidence that
Kahneman-Tversky prospect theory confers OOS safety on RL trading policies.
\normalsize

\newpage
\tableofcontents
\newpage

% ─── 1. Introduction ─────────────────────────────────────────────────────────
\section{서론}
\label{sec:intro}

BTC 등 암호화폐 시장의 시간봉 가격 변동은 장기 추세와 무관한 빈번한 미세 진동
(micro-oscillation)이 본질적 특성으로, 이 변동성 자체를 수익원으로 삼는
\emph{그리드 트레이딩}(grid trading)은 시장 조성(market making)
\citep{avellaneda2008market}의 단순화된 변형이다. 고정된 가격 격자 위에 매수/매도
주문을 정기적으로 배치하는 단순한 규칙 기반 전략부터, ATR(Average True Range)
\citep{wilder1978atr} 비례로 격자 폭을 동적 조정하는 변동성 적응형 변형까지 다양한
구현이 존재한다. 본 연구는 후자, 즉 ATR 비례 그리드 트레이딩의 \emph{격자 폭과
체결 가격을 강화학습 정책으로 결정}하는 시스템을 다룬다.

\subsection{동기와 연구 질문}
\label{sec:intro_motivation}

본 시스템의 초기 가설(2026년 1\textendash 3월 작업)은 ``PPO 강화학습이 동적
변동성 적응을 학습하여 고정 그리드 베이스라인을 초과한다''였다. 그러나 사후 분석
결과 학습된 정책이 사실상 상수 액션 $[1.0, 0.0]$으로 수렴함이 확인되었고
($\S$\ref{sec:phase1_negative} 상세), Bayesian으로 최적화된 ATR 계수 자체가 성능의
주요 원동력임이 드러났다. 이 negative finding은 \emph{대칭 보상(symmetric reward)
$+$ ATR 비례 공식의 조합이 RL 정책의 자유도를 본질적으로 흡수}한다는 가설로 이어졌고,
본 논문의 새 연구 질문을 동기화한다:

\begin{quote}
\textbf{BTC 그리드 트레이딩에서 reward 함수의 설계가 RL 정책의 행동 패턴과 일반화
성능에 어떤 영향을 미치는가? 특히, 손실 비대칭(asymmetric, prospect-theoretic) 또는
위험 조정 수익(DSR) 형태의 reward 정식화가 단일 metric 우위 또는 분포 외(out-of-sample)
강건성 측면에서 의미 있는 차이를 만드는가?}
\end{quote}

\subsection{기여 사항}
\label{sec:intro_contribution}

본 논문의 학술적 기여는 다음 네 가지이다.

\begin{enumerate}[leftmargin=*]
\item \textbf{Reward 변형의 통계적으로 엄밀한 비교 프레임워크}: 4가지 reward
함수(\texttt{sym}, \texttt{asym}, \texttt{dsr}, \texttt{pt})를 동일 환경에서
10 시드 $\times$ 100만 스텝으로 학습하고, Cohen's $d$, 부트스트랩 P(A$>$B), IQM,
5\% CVaR을 사용하여 페어와이즈 통계 검정을 수행한다.
\item \textbf{시나리오 D --- Pareto 프론티어 발견}: 단순 ``X reward가 best''의
단일 winner 결론 대신, 4 변형이 \{\texttt{sym}, \texttt{dsr}\}와
\{\texttt{asym}, \texttt{pt}\}의 두 클러스터로 통계적으로 분리되어 Sharpe--MDD
평면에서 Pareto-유사 프론티어를 형성함을 보이며($\S$\ref{sec:pareto}), 이는 사전에
설정한 시나리오 분기(낙관 A/중립 B/비관 C) 어디에도 해당하지 않는 새로운 결과 형태이다.
\item \textbf{Reward $\to$ 행동 $\to$ 결과의 인과 사슬 정량}: 변형 간 정책 행동의
차이를 1.04M step trajectory 위에서 정책 거리 행렬(within-cluster 0.129 vs.\
across-cluster 0.286), regime별 거래/holding 통계, 행동 분포 비교로 정량화하여
``reward 형식의 손실 비대칭 $\to$ 거래 빈도 감소 $\to$ 보수적 클러스터''의 메커니즘을
확인한다($\S$\ref{sec:mechanism}).
\item \textbf{평가 환경별 winner reversal과 Prospect Theory의 OOS 강건성}:
Single-split Val(\texttt{sym} 1위) $\to$ CPCV multi-split(\texttt{dsr} 1위) $\to$
봉인 해제된 Test(\texttt{pt} 1위, 두 소스 일관 p$<$0.002)의 세 환경에서 winner가
완전 reversal한다($\S$\ref{sec:robustness}). \texttt{pt}의 OOS 강건성 메커니즘은
prospect theory의 손실 회피 $\lambda=3.30$과 오목 효용 $\alpha=0.68$이 정책의
holding을 매우 짧게(평균 1.4시간) 학습시켜 미관측 강세장의 매도 타이밍 위험을
회피하는 결과임을 trajectory 분석으로 입증한다. 이는 Kahneman과 Tversky(1979)의
이론이 강화학습 정책에 부여하는 분포 외 안전성의 첫 정량 확인이다.
\end{enumerate}

\subsection{본 논문 구성}
\label{sec:intro_structure}

\S\ref{sec:related}는 RL 트레이딩, reward 설계 이론, 평가 방법론의 선행 연구를
정리한다. \S\ref{sec:method}는 환경 설계(MDP, ATR 공식, 4 reward 변형)와 PPO 학습
설정을 기술한다. \S\ref{sec:phase1_negative}는 본 연구의 동기인 Phase 1 negative
finding(``RL $=$ Fixed $[1.0, 0.0]$'')을 간략 재인용한다. \S\ref{sec:pareto}는
Val에서의 두 클러스터 발견(시나리오 D)을, \S\ref{sec:mechanism}은 메커니즘 정량을
다룬다. \S\ref{sec:robustness}은 slippage 견고성(\S\ref{sec:slippage}), CPCV
다중 split(\S\ref{sec:cpcv}), 그리고 본 논문의 핵심인 Test OOS와
prospect theory의 강건성 메커니즘(\S\ref{sec:oos})을 다룬다. \S\ref{sec:discussion}은
분포 이동의 정량과 한계를 논의하며, \S\ref{sec:conclusion}에서 마무리한다.

% ─── 2. Related Work ─────────────────────────────────────────────────────────
\section{관련 연구}
\label{sec:related}

본 연구는 (i) 시장 조성과 그리드 트레이딩, (ii) 강화학습 기반 트레이딩,
(iii) reward 설계 이론, (iv) 트레이딩 정책의 평가 방법론의 네 분야 교차점에
위치한다. 이 절은 각 분야의 핵심 선행 연구를 정리하고, 본 논문이 차지하는
좌표를 명시한다.

\subsection{시장 조성과 그리드 트레이딩}
\label{sec:related_grid}

그리드 트레이딩의 학술적 조상은 \citet{avellaneda2008market}의 시장 조성
(market making) 모델이다. 이 연구는 위험 회피적 시장 조성자가 재고 위험(inventory
risk)에 직면할 때, 중심 가격(mid-price)을 중심으로 매수/매도 호가를
\emph{대칭적이지 않게} 배치하는 최적 정책의 폐형 해(closed-form solution)를
유도한다. 핵심 통찰은 (a) 호가의 폭(spread)이 변동성과 시간 잔여량의 함수로
결정되며, (b) 재고 편향(inventory skew)이 호가 중심을 비대칭적으로 이동시킨다는
것이다. 그리드 트레이딩은 이 모델의 단순화된 실용 변형으로, 고정 간격 격자에
주문을 미리 배치하여 미세 변동성을 수익화한다.

격자 폭을 시장 변동성에 적응시키는 표준 방법은 \citet{wilder1978atr}의
ATR(Average True Range, 14일 이동평균 윈도우 표준)을 사용하는 것이다. ATR은
실제 변동 폭(true range)의 평활 평균으로 변동성 클러스터링(volatility clustering)
구간에서 격자 폭을 자동 확장하는 단순하고도 견고한 규칙을 제공한다.
ATR 비례 그리드 ($\Delta = c \cdot \mathrm{ATR}$, $c$ = 비례 계수)는 실무에서
널리 사용되며, 본 연구의 베이스라인이기도 하다.

이 분야의 기존 학술 연구는 대부분 \emph{규칙 기반(rule-based)} 정책의 분석에
머물러 있으며, 격자 폭 결정을 데이터로부터 학습하는 시도는 제한적이다. 본 논문은
ATR 비례 공식의 \emph{비례 계수와 격자 분할 수를 강화학습 정책으로 동적 결정}하는
구조를 채택하면서도, 동시에 reward 함수 설계가 학습된 정책의 행동에 어떤 영향을
미치는지 정량적으로 분석한다는 점에서 차별성을 가진다.

\subsection{강화학습 기반 트레이딩}
\label{sec:related_drl}

암호화폐 시장에 대한 DRL 트레이딩의 표준 베이스라인은 \citet{zhang2020drl}이
제시한 PPO/A2C/DQN 비교 프레임워크이다. 이 연구는 가격 시계열을 입력으로
직접 받는 PPO 정책이 단순 buy-and-hold(B\&H)와 무빙 평균 교차 전략을 일관적으로
초과하는 결과를 보였다. \citet{sun2023rl}의 ACM Transactions 서베이와
\citet{liu2025review}의 최근 리뷰는 이 분야의 두 가지 미해결 과제로
\emph{(a)~정책 행동의 해석 가능성, (b)~시장 레짐 간 일반화}를 지적한다.
본 논문은 이 두 과제를 직접 다룬다: §\ref{sec:mechanism}에서 reward 형식에서
행동 차이로 가는 인과 사슬을 정량화하며, §\ref{sec:robustness}에서 단일 분할
검증으로 잡히지 않는 분포 외(OOS) 일반화 현상을 보고한다.

직접 선행 연구는 네 편이다. \citet{liu2021bitcoin}은 LSTM 가격 예측의 출력을
PPO 정책의 state로 입력하는 2단계 구조를 제안했다. 이 접근의 한계는 예측 오류가
정책 결정에 그대로 전파된다는 점이며, 본 논문은 가격 예측 단계를 \emph{완전 제거}한
변동성 적응형 그리드 구조로 이 한계를 우회한다. \citet{yasin2024rl}은 DQN/PPO/A2C
세 알고리즘과 20개 기술지표를 비교하나, action 공간이 이산 buy/sell이고
holding 행동이 없는 한계가 있다. 본 논문은 연속 2차원 action(공격성과 이익 목표)을
채택해 그리드의 미세 조정을 가능하게 한다. \citet{pham2025dqn}은 사전 정의된
5개 전략 중 하나를 DQN으로 선택하는 메타-전략 구조이나, 전략 자체는 학습되지
않는다. \citet{bandarupalli2025risk}는 PPO에 변동성 패널티와 거래 비용 항을
추가하나, 단일 reward 변형의 hyperparameter 조정 수준에 머물러 본 논문이
다루는 reward 형식 자체의 비교는 포함하지 않는다.

특히 중요한 선행 연구는 \citet{gort2022backtest}이다. 이 논문은 암호화폐 DRL
트레이딩에서 \emph{단일 train/test 분할의 결과가 OOS 일반화를 보장하지 않는다}는
실증을 제시하며, CPCV(Combinatorial Purged Cross-Validation)와 DSR(Deflated
Sharpe Ratio)을 권장한다. 본 논문은 Gort et al.의 평가 권장사항을 채택하여 Val
(2021--2023) / CPCV 6-fold (Train 분할 내부) / Test (2024-- 봉인 해제) 세 환경을
동시에 사용한다. 본 논문의 \emph{winner reversal} 발견(Val sym → CPCV dsr → Test
pt)은 Gort et al.의 우려가 reward 변형 선택의 차원에서도 동일하게 나타남을
정량 확인한 것이다.

본 논문의 위치를 한 줄로 요약하면: 기존 DRL 트레이딩 문헌이 \emph{단일 reward
함수의 알고리즘/state 설계 비교}에 집중한 반면, 본 논문은 \emph{같은 알고리즘
(PPO)과 같은 환경 위에서 reward 함수 자체의 4가지 변형을 통제 비교}하며,
그 효과를 in-sample / multi-split / OOS 세 환경에서 분리 측정한다.

\subsection{Reward 설계 이론}
\label{sec:related_reward}

본 논문이 비교하는 4가지 reward 변형은 강화학습과 행동경제학의 세 가지 이론적
계보에서 도출된다. (i) 대칭(\texttt{sym})은 표준 PnL 수익률 reward로 비교 베이스라인
역할을 한다. (ii) 비대칭(\texttt{asym})은 손실에 양의 가중치 $\beta>1$을 부여하는
변형으로, 위험 회피 효용 함수의 단순 1차 근사이다. (iii) 미분 샤프 비율(\texttt{dsr})은
\citet{moody2001dsr}의 정식이다. (iv) 프로스펙트 이론 보상(\texttt{pt})은
\citet{kahneman1979prospect}과 \citet{tversky1992cumulative}의 효용 함수
$v(x) = \mathrm{sgn}(x) |x|^\alpha$ ($x \geq 0$일 때) 또는
$-\lambda |x|^\alpha$ ($x < 0$일 때)를 단계 reward로 직접 사용한다.

\citet{moody2001dsr}의 DSR은 샤프 비율의 \emph{온라인 추정치 변화량}을 step reward로
사용하는 정식으로, 지수 가중 이동 평균(EWMA, decay rate $\eta$)으로 1차와 2차
모멘트를 추적한다. 이 정식은 강화학습 정책 학습에 두 가지 의미를 갖는다: (a) 단순
PnL 대비 위험 조정 수익을 직접 최적화한다, (b) sliding window의 메모리 구조가
정책에 \emph{단일 step 이상의 시간 의존성}을 부여한다. 본 논문 §\ref{sec:mechanism}은
이 메모리 구조의 결과로 \texttt{dsr} 변형의 정책이 다른 변형 대비 2--6배 긴
holding 시간을 학습함을 정량 확인하며, 이는 §\ref{sec:robustness}의 in-sample
다중 분할에서의 우위와 §\ref{sec:oos}의 OOS 실패로 이어지는 인과의 출발점이다.

\citet{kahneman1979prospect}의 프로스펙트 이론은 인간 의사 결정의 위험 행동을
설명하는 행동경제학 기초이며, 특히 (a) 손실 회피 $\lambda \approx 2.25$
(\citealt{tversky1992cumulative}의 표준값), (b) 효용 함수의 오목성 $\alpha \approx 0.88$
이 핵심 모수이다. 이 이론은 인간 트레이더의 행동(예: ``손실에 매달리고 이익은 빨리
실현'')을 설명하는 데는 광범위하게 사용되어 왔으나, \emph{강화학습 정책의
reward로 직접 사용}한 학술 연구는 매우 드물다. 본 논문은 \citet{kahneman1979prospect}의
정식을 RL reward로 직접 구현하고 ($\alpha=0.68$, $\lambda=3.30$, Optuna로 튜닝),
이 reward로 학습된 정책이 분포 외 환경에서 보이는 강건성을 정량적으로
입증한다는 점에서, 본 논문이 알기로는 첫 사례에 해당한다.

\citet{ng1999reward}의 정책 불변 reward shaping(potential-based shaping) 정리는
``보상의 단조 변환은 정책 학습 결과에 영향을 미치지 않는다''는 결론을 도출한다.
얼핏 본 논문의 4 reward 변형 비교가 이 정리와 충돌하는 듯 보이지만, 그렇지
않다: \texttt{sym} $\to$ \texttt{asym}/\texttt{pt}의 변환은 음의 영역에 비대칭
스케일을 부여하는 \emph{비-potential-based} 변환이며, \texttt{dsr}은 단순 step
reward 변환이 아닌 \emph{state-extended} reward(즉, sliding window 모멘트가
사실상 state의 일부)이다. 따라서 본 논문의 reward 변형 비교는 Ng et al.의
정리의 적용 범위 밖에 있으며, 실제로 변형 간 정책 행동이 통계적으로 유의하게
차이남을 §\ref{sec:mechanism}에서 확인한다.

\subsection{평가 방법론}
\label{sec:related_eval}

강화학습 정책의 평가에서 가장 큰 함정은 (a) 백테스트 과적합과 (b) 결과의
시드 의존성이다. \citet{henderson2018rlreproducibility}는 동일 알고리즘이라도
시드 차이만으로 결과가 크게 변동함을 실증하고, 5--10 시드의 통계적 비교, IQM
(Interquartile Mean), 부트스트랩 신뢰구간을 표준 권장사항으로 제시했다. 본
논문은 변형당 10 시드(Val/exp033), 100 모델(Test)을 사용하며 페어와이즈
Cohen's $d$와 부트스트랩 $\mathrm{P}(A>B)$를 모든 비교에 적용한다.

\citet{lopezdeprado2018advances} (8장)의 \emph{Combinatorial Purged
Cross-Validation}(CPCV)은 시계열 데이터에서 train/test 누설을 방지하면서
$\binom{N}{k}$개의 분할 조합을 활용하여 단일 분할 결과의 일반화 가능성을
강화하는 방법이다. 본 논문은 6-fold CPCV(15 paths)를 §\ref{sec:cpcv}에서 채택한다.

\citet{lopezdeprado2014dsr}의 \emph{Deflated Sharpe Ratio}(DSR\textsuperscript{*})는
백테스트 시도 횟수와 분포의 첨도/왜도를 보정한 샤프 비율의 통계적 유의성
검정이다. 본 논문은 CPCV 결과의 다중 비교에 Bonferroni 보정 4-way + DSR\textsuperscript{*}
검정을 함께 사용하며 (§\ref{sec:cpcv}), 모든 reward 변형이 ATR 대비 $p<0.001$의
유의성을 만족함을 확인한다.

본 논문의 평가 방법론 측면 차별점은 \emph{세 환경(Val / CPCV / Test)을 동시에
사용하여 winner reversal을 발견}한 점이다. 단일 분할 또는 CPCV 단독을 사용하는
기존 연구는 OOS 일반화 실패를 발견할 수 없으며, OOS만 사용하는 연구는 \texttt{pt}
변형의 in-sample 보수성을 발견할 수 없다. 세 환경의 동시 적용은
\citet{henderson2018rlreproducibility}와 \citet{gort2022backtest}이 권장하는
복수 평가 방식을 \emph{reward 변형 비교라는 단일 차원에 일관 적용}한 점에서
방법론적 기여를 가진다.

% ─── 3. Method ───────────────────────────────────────────────────────────────
\section{방법론}
\label{sec:method}

\textit{[TBD: 추후 작성. 서브 섹션:
\S3.1 MDP 설계 (State 7D, Action 2D, ATR 비례 공식),
\S3.2 Trading 환경 구체 (사이클, fee, slippage),
\S3.3 4 reward 변형 정의 + Optuna 튜닝 (exp032a),
\S3.4 PPO 안정화 패키지 (target\_kl, ent annealing),
\S3.5 평가 방법론 (Val/CPCV/Test).]}

% ─── 4. Phase 1 Negative Finding ─────────────────────────────────────────────
\section{Phase 1: 대칭 보상 하의 정책 포화}
\label{sec:phase1_negative}

\textit{[TBD: 추후 작성. exp020--022의 ``RL = Fixed [1.0, 0.0]'' 발견.
대칭 보상 + ATR 비례 공식 조합의 RL 자유도 흡수.
본 논문의 reward 정식화 동기.]}

% ─── 5. Main Finding: Pareto Frontier ────────────────────────────────────────
\section{Pareto 프론티어와 시나리오 D}
\label{sec:pareto}

\textit{[TBD: 추후 작성. exp032b 결과:
4 variants $\times$ 10 seeds 통계, Pareto scatter,
두 클러스터 (aggressive/conservative), 가설 H1\textendash H3 점검,
시나리오 D 확정.]}

% ─── 6. Mechanism Quantification ─────────────────────────────────────────────
\section{메커니즘 정량화}
\label{sec:mechanism}

\textit{[TBD: 추후 작성. exp032c 결과:
Policy distance 행렬, behavior per regime,
counterfactual action map, action distribution.
인과 사슬: reward 형식 $\to$ 거래 빈도 $\to$ 클러스터.]}

% ─── 7. Robustness Trio ──────────────────────────────────────────────────────
\section{강건성 분석}
\label{sec:robustness}

\subsection{Slippage 강건성}
\label{sec:slippage}

\textit{[TBD: 추후 작성. exp033: 4 variants $\times$ 10 seeds with slippage 0.02\%.
Cluster preservation ratio 2.19$\times$. ATR-with-slippage fair 비교 결과
(Phase 16a).]}

\subsection{CPCV 다중 split 평가}
\label{sec:cpcv}

\textit{[TBD: 추후 작성. exp034: 6-fold CPCV 15 paths.
DSR p$<$0.001 (Bonferroni 4-way 보정), \texttt{dsr} variant reversal 1위
(1.413, 5\% CVaR 0.890), Cluster ratio 3.55$\times$.]}

\subsection{분포 외 평가와 Prospect Theory의 강건성}
\label{sec:oos}

\textit{[TBD: 추후 작성. 본 논문의 핵심 챕터.
exp035: 100 RL models + ATR baseline + Buy-and-Hold Test.
\texttt{pt}의 Test 1위 (두 소스 p$<$0.002), \texttt{dsr}의 Test 꼴찌.
Phase 16d mechanism: holding 시간 분포 (DSR 7-day max vs.\ pt 6h max).
\texttt{pt}의 분포 외 강건성의 인과 사슬.]}

% ─── 8. Discussion ───────────────────────────────────────────────────────────
\section{논의}
\label{sec:discussion}

\textit{[TBD: 추후 작성. 서브 섹션:
\S8.1 Val--Test 분포 이동의 정량 (KS p$<10^{-10}$, $-27\%$ variance),
\S8.2 exp027\_rl 사전 증거(Env-v3 asym 1.955)의 환경 의존성 정직 인정,
\S8.3 단일 metric winner의 한계,
\S8.4 한계: BTC 단일 자산, Test 1 regime, single slippage level.]}

% ─── 9. Conclusion ───────────────────────────────────────────────────────────
\section{결론}
\label{sec:conclusion}

\textit{[TBD: 추후 작성. 4개 메인 메시지:
(1) Pareto 프론티어 trade-off,
(2) 평가 환경별 winner reversal,
(3) Prospect Theory의 OOS 강건성,
(4) In-sample 우위 $\neq$ OOS 강건성.
Future work: multi-asset, regime, DR, meta-RL.]}

% ─── References ──────────────────────────────────────────────────────────────
\bibliographystyle{plainnat}
\bibliography{references}

\end{document}
